\documentclass[12pt, a4paper, oneside, notitlepage]{report}
\usepackage[T1]{fontenc}
\usepackage[utf8]{inputenc}
\usepackage{lmodern}
\usepackage{amsmath, amssymb, amsthm}
\usepackage{graphicx}
\usepackage{hyperref}
\usepackage{cleveref}
\usepackage{cite}
\usepackage{url}
\usepackage{booktabs}
\usepackage{microtype}
\usepackage{listings}
\usepackage{color}
\usepackage{xcolor}
\usepackage{geometry}
\geometry{margin=1.2in}

% =============================================================================
% Macros
% =============================================================================
\definecolor{codegreen}{rgb}{0,0.6,0}
\definecolor{codegray}{rgb}{0.5,0.5,0.5}
\definecolor{codepurple}{rgb}{0.58,0,0.82}
\lstset{basicstyle=\ttfamily\small, numbers=left, stepnumber=1,
  showstringspaces=false, breaklines=true, frame=single,
  backgroundcolor=\color{white}, keywordstyle=\color{blue},
  commentstyle=\color{codegreen}, stringstyle=\color{codepurple}}

\theoremstyle{plain}
\newtheorem{definition}{Definition}[chapter]
\newtheorem{assumption}{Assumption}[chapter]
\newtheorem{example}{Example}[chapter]

\DeclareMathOperator*{\argmax}{arg\,max}
\DeclareMathOperator*{\argmin}{arg\,min}
\DeclareMathOperator{\tr}{tr}
\DeclareMathOperator{\id}{id}
\DeclareMathOperator{\Real}{Re}
\DeclareMathOperator{\Imag}{Im}
\DeclareMathOperator{\supp}{supp}

% =============================================================================
% Title Page
% =============================================================================
\begin{document}
\pagenumbering{roman}

\begin{titlepage}
  \centering
  \vspace*{2cm}
  {\LARGE\bfseries
    Quantum Reservoir Computing for EV Charging Load Forecasting\\
    with Quantum Hidden Markov Models and Eligibility Traces\par}
  \vspace{1.5cm}
  {\large\bfseries Riza Syah\par}
  \vspace{0.5cm}
  {\large Department of Quantum Computing\par}
  \vspace{2cm}
  {\large A paper submitted in partial fulfillment\\
    of the requirements for the degree of\\
    Master of Science in Quantum Artificial Intelligence\par}
  \vspace{2cm}
  {\large\.today\par}
  \vspace{1.5cm}
  \textbf{Abstract:}
  \parbox{0.85\textwidth}{\small
    We present QRC-EV, a complete Quantum Reservoir Computing framework for
    forecasting electric vehicle charging load. The framework combines a
    stateful GPU quantum reservoir (CUDA Quantum on NVIDIA hardware via
    Amazon Braket) with a Quantum Hidden Markov Model (QHMM) whose
    transition dynamics are learned via Online Maximum Likelihood
    Estimation (OMLE) subject to complete positivity and trace
    preservation (CPTP) constraints. The QHMM is represented as an
    Observable Operator Model (OOM), enabling exact trajectory probability
    computation via a forward-backward algorithm. We extend the framework
    with TD($\lambda$) eligibility traces for temporal credit assignment,
    allowing the model to learn from ongoing interaction without
    backpropagation through the quantum processor. Comprehensive
    benchmarks on synthetic (sinusoidal, Mackey-Glass, NARMA-10) and
    real-world EV charging datasets (ACN-Data, Palo Alto) demonstrate
    competitive forecasting performance compared to classical baselines
    including Echo State Networks and LSTMs.}
  \clearpage
\end{titlepage}

% =============================================================================
% Acknowledgments
% =============================================================================
\chapter*{Acknowledgments}
\addcontentsline{toc}{chapter}{Acknowledgments}
I acknowledge the support of the NVIDIA CUDA Quantum team and Amazon Web
Services for providing access to quantum-classical computing infrastructure.
Special thanks to the members of the Quantum AI reading group for valuable
feedback on early drafts.

% =============================================================================
% Table of Contents
% =============================================================================
\tableofcontents
\listoffigures
\listoftables

% =============================================================================
% Main text
% =============================================================================
\pagenumbering{arabic}
% =============================================================================
% Chapter 1 — Introduction
% =============================================================================

\chapter{Introduction}
\label{chap:intro}

% =============================================================================
% §1.1  Context and Motivation
% =============================================================================

\section{Context and Motivation}
\label{sec:intro:context}

The global transition toward sustainable transportation has produced an
unprecedented growth in Electric Vehicle (EV) adoption. Managing the
resulting charging infrastructure presents a \emph{time-series forecasting
challenge} of considerable practical and economic magnitude: grid operators
must balance load in real time, charging networks must schedule capacity,
and energy providers must procure power days in advance. Accurate
forecasting of EV charging demand—across hourly, daily, and weekly
horizons—is therefore a critical enabling technology
\citep{acaroglu2022acn, chen2020ev}.

Classical machine learning approaches to this problem include recurrent
neural networks (RNNs), Long Short-Term Memory networks (LSTMs)
\citep{hochreiter1997lstm}, and Transformer-based models such as the
Temporal Fusion Transformer (TFT) \citep{lim2019tft}. These methods achieve
strong performance on short-term prediction but require extensive
training data, large parameter counts, and significant computational
resources for inference at scale. Echo State Networks (ESNs)
\citep{jaeger2004harnessing}—a form of reservoir computing—offer a
lightweight alternative in which only the readout layer is trained,
leaving the reservoir dynamics fixed. ESNs exploit the echo state
property to reliably encode temporal dependencies without gradient-based
learning.

Despite their appeal, ESNs operate on classical hardware, and their
representational capacity is fundamentally bounded by the classical
dynamics of their reservoir. A natural question is whether quantum
systems, whose Hilbert-space dimension grows exponentially with the
number of qubits, can provide richer temporal representations at
comparable or lower physical resource cost.

% =============================================================================
% §1.2  Quantum Reservoir Computing
% =============================================================================

\section{Quantum Reservoir Computing}
\label{sec:intro:qrc}

Quantum Reservoir Computing (QRC) \citep{jaeger2004harnessing,
fujii2017keep, appert2023quantum} proposes to replace the classical
reservoir with a quantum system: the input drives the evolution of a
quantum register, and the register's state at each timestep is measured to
produce a feature vector for a classical readout. If the quantum system
exhibits the \emph{quantum echo state property}—analogous to the echo
state property—its dynamics act as a universal black-box model for
temporal computation.

The practical appeal of QRC rests on two pillars. First, the quantum
system is \emph{not trained}; only the measurement operators and the
classical readout are optimized. This sidesteps the classical
barriers to training quantum neural networks in the NISQ era. Second,
quantum dynamics—particularly those near the edge of quantum chaos—can
generate entanglement and complex correlations that are believed to exceed
what classically simulable systems can produce efficiently
\citep{arute2019quantum, boixo2018characterizing}. This motivated a
series of experimental demonstrations showing quantum advantage for
specific temporal tasks
\citep{ Dumitrescu2018, Chen2020, Zheng2021, Ollitrault2023}.

Nevertheless, the practical advantage of QRC on real-world forecasting
tasks—particularly in the high-stakes, noisy regime of energy
infrastructure—remains an open question. Published benchmarks are
dominated by synthetic tasks (Mackey-Glass, NARMA) and small-scale
laboratory settings, leaving a significant gap between controlled
experiments and deployable forecasting systems
\citep{gosavic2024quantum, gonzalez2024benchmark}.

% =============================================================================
% §1.3  This Work
% =============================================================================

\section{This Work: QRC-EV Forecasting System}
\label{sec:intro:this_work}

This paper presents a complete \textbf{Quantum Reservoir Computing
framework for EV charging load forecasting} (QRC-EV) that bridges this
gap. Our contributions are organized along three dimensions:

\subsection{Quantum-Classical Architecture}
We implement a \textbf{stateful GPU quantum reservoir} using NVIDIA's
CUDA Quantum platform \citep{nvidiacudaq} integrated with Amazon Braket
\citep{amazonbraket}. The reservoir operates on up to 30 qubits with
input-scaled parametric gates, maintaining numerical stability across
long input sequences. The reservoir state is evolved via a
Lindblad-style dissipative channel, producing a feature map that is
subsequently processed by a classical Ridge regression readout. This
design runs entirely on cloud GPU infrastructure without the need for
physical quantum hardware access, enabling reproducible large-scale
experiments.

\subsection{Quantum Hidden Markov Model with Eligibility Traces}
Purely reservoir-based approaches treat the quantum state as a passive
temporal feature extractor. We instead pair the reservoir with a
\textbf{Quantum Hidden Markov Model (QHMM)} whose parameters are learned
via \textbf{Online Maximum Likelihood Estimation (OMLE)} subject to
complete positivity and trace preservation (CPTP) constraints. The QHMM
is represented as an \textbf{Observable Operator Model (OOM)} whose
transition structure is learned from trajectory data using a
forward-backward algorithm for state smoothing, analogous to the
Baum-Welch algorithm for classical HMMs.

Critically, we extend the QHMM-OOM framework with
\textbf{TD($\lambda$) eligibility traces} for temporal credit assignment.
The forward trace accumulates the forward-message state history;
the backward trace propagates TD errors; their product enables
efficient parameter updates without full backpropagation through time.
This allows the QHMM to learn from ongoing interaction, making it
suitable for online learning settings.

\subsection{Comprehensive Benchmarks}
We evaluate QRC-EV on four datasets spanning synthetic benchmarks
(sinusoidal, Mackey-Glass, NARMA-10) and real-world EV charging data
(ACN-Data, Palo Alto network). We compare against classical baselines
(Ridge regression, ESN, LSTM) and ablate the contribution of each
component (reservoir size, OOM smoothing, eligibility traces). Our
results demonstrate that the QHMM-OMLE pipeline achieves competitive
performance on multi-step forecasting, particularly at horizons where
the classical reservoir provides an informative prior.

% =============================================================================
% §1.4  Paper Structure
% =============================================================================

\section{Paper Structure}
\label{sec:intro:structure}

The paper is organized as follows. Chapter~\ref{chap:background}
reviews the necessary background: reservoir computing, quantum
channels and the Choi-Jamiołkowski representation, the HMM-to-OOM
connection, and the forward-backward algorithm. Chapter~\ref{chap:methods}
presents our methods: the GPU reservoir architecture, the QHMM-OOM
formalism with OMLE, the TD($\lambda$) trace mechanism, and the
optimistic planning algorithm for prediction.
Chapter~\ref{chap:experiments} describes the experimental setup,
datasets, baselines, and evaluation protocol.
Chapter~\ref{chap:results} presents and discusses the results.
Chapter~\ref{chap:conclusion} concludes with implications and open
questions.

% =============================================================================
% §1.5  Reproducibility Statement
% =============================================================================

\section*{Reproducibility Statement}
\label{sec:intro:reproducibility}

All code is publicly available at
\url{https://github.com/rezacute/medseg3d-research}. Experiments use
CUDA Quantum 0.6+ with Amazon Braket and can be run on any NVIDIA
GPU with compute capability $\geq 8.0$. Hyperparameters and random
seeds are documented in Appendix~\ref{app:hyperparameters}.

% =============================================================================
% Chapter 2 — Background and Mathematical Foundations
% =============================================================================

\chapter{Literature and Mathematical Foundations}
\label{chap:background}

This chapter establishes the mathematical and conceptual foundations
necessary for the QRC-EV framework. We begin with reservoir computing
(\S\ref{sec:rc}) and its quantum counterpart (\S\ref{sec:qrc_formalism}),
then develop the theory of quantum channels and the Choi-Jamiołkowski
representation (\S\ref{sec:quantum_channels}), the HMM-to-OOM
connection (\S\ref{sec:hmm_oom}), the forward-backward algorithm
(\S\ref{sec:forward_backward}), and finally TD($\lambda$) eligibility
traces (\S\ref{sec:td_lambda}).

% =============================================================================
% §2.1  Classical Reservoir Computing
% =============================================================================

\section{Classical Reservoir Computing and Echo State Networks}
\label{sec:rc}

Reservoir Computing (RC) refers to a family of temporal machine learning
frameworks in which a \emph{fixed} dynamical system—the \emph{reservoir}—
maps an input time series to a high-dimensional feature space; a
subsequent \emph{readout} layer is the only component that is trained
\citep{jaeger2004harnessing, jaeger2007echo, schrauwen2007rc}. The key
property enabling learning without backpropagation through the reservoir
is the \textbf{echo state property} (ESP):

\begin{definition}[Echo State Property]
\index{echo state property}
A reservoir with state update $r_{t+1} = f(W_{\mathrm{in}} u_{t+1} + W r_t)$
has the echo state property if, for any two input sequences
$\{u_t\}$ and $\{u'_t\}$ that differ only in the distant past,
the resulting reservoir states converge to each other as $t \to \infty$,
independently of the initial reservoir state $r_0$.
\end{definition}

When the ESP holds, the reservoir acts as a \emph{universal temporal
basis}: any bounded function of the input history can be approximated
by a linear combination of reservoir features. Training reduces to solving
a linear regression problem.

An Echo State Network (ESN) instantiates this with a classical recurrent
network:
\begin{equation}
r_{t+1} = (1 - \alpha) \, r_t + \alpha \, \tanh\!\bigl( W_{\mathrm{in}} u_{t+1}
  + W r_t + W_{\mathrm{bias}} \bigr),
\qquad
y_t = W_{\mathrm{out}} \, r_t,
\label{eq:esn_update}
\end{equation}
where $\alpha \in (0,1]$ is the leaking rate, $W_{\mathrm{in}} \in
\mathbb{R}^{N \times d_{\mathrm{in}}}$ maps the $d_{\mathrm{in}}$-
dimensional input to the $N$-dimensional reservoir, and $W \in
\mathbb{R}^{N \times N}$ is the reservoir weight matrix with spectral
radius $\rho(W) < 1$ to ensure the ESP. Only $W_{\mathrm{out}}}$ is
trained, typically via Ridge regression.

ESNs have demonstrated competitive performance on a range of temporal
learning tasks including system identification, speech recognition, and
chaotic time-series forecasting \citep{jaeger2004harnessing,
verstraeten2007esn}. Their main limitation is the representational
capacity of the reservoir, which is fundamentally bounded by the classical
dynamics of the weight matrix.

% =============================================================================
% §2.2  Quantum Reservoir Computing
% =============================================================================

\section{Quantum Reservoir Computing}
\label{sec:qrc_formalism}

Quantum Reservoir Computing (QRC) replaces the classical reservoir with a
quantum many-body system \citep{fujii2017keep, appert2023quantum,
ghojeh2024quantum}. Consider an $n$-qubit quantum register initialized in
a fixed state $\rho_0$. An input $u_t \in \mathbb{R}^{d_{\mathrm{in}}}$ is
encoded as a parameterized unitary $U(u_t)$ acting on the register:

\begin{equation}
\rho_{t+1} = \mathcal{E}\!\bigl( U(u_{t+1}) \, \rho_t \, U(u_{t+1})^\dagger \bigr),
\label{eq:qrc_update}
\end{equation}
where $\mathcal{E}$ is a dissipative channel (typically a depolarizing or
amplitude-damping channel) that ensures the register relaxes toward a
stationary state and prevents indefinite purity growth. The reservoir
output at time $t$ is obtained by measuring a fixed set of observables
$\{O_j\}$, producing a feature vector
\begin{equation}
\phi_t = \bigl( \operatorname{tr}[O_1 \rho_t],
               \ldots, \operatorname{tr}[O_m \rho_t] \bigr)^\top
\in \mathbb{R}^m.
\end{equation}
The readout $y_t = W_{\mathrm{out}} \phi_t$ is trained exactly as in
the classical ESN case.

The quantum echo state property (QESP) is defined analogously to the
classical case \citep{fujii2017keep}: the state $\rho_t$ should become
insensitive to the initial state $\rho_0$ as $t$ increases, provided the
input is bounded. For a system near a unique fixed point of the channel
$\mathcal{E}$, the QESP holds automatically; for systems with multiple
stationary states the condition becomes non-trivial.

The advantage of QRC over classical RC rests on the exponential
dimensionality of the Hilbert space ($\dim \mathcal{H} = 2^n$) and the
capacity of quantum dynamics to generate entanglement. However, the
accessible information in $\phi_t$ is ultimately limited by the
measurement observables and by decoherence. The practical question is
whether the \emph{effective} dimensionality of the quantum feature space
exploited by the readout exceeds what a classical reservoir of comparable
size can achieve.

Recent experimental work has demonstrated QRC on physical quantum
hardware including photonic systems \citep{Dumitrescu2018}, trapped ions
\citep{Chen2020}, and superconducting qubits \citep{Zheng2021}. These
proofs-of-concept are encouraging but typically evaluate on low-dimensional
synthetic tasks. The benchmark gap between controlled experiments and
real-world forecasting remains a significant obstacle
\citep{gosavic2024quantum, gonzalez2024benchmark}.

% =============================================================================
% §2.3  Quantum Channels and the Choi-Jamiołkowski Representation
% =============================================================================

\section{Quantum Channels and the Choi-Jamiołkowski Representation}
\label{sec:quantum_channels}

A \textbf{quantum channel} $\Phi: \mathcal{B}(\mathcal{H}_{\mathrm{in}})
\to \mathcal{B}(\mathcal{H}_{\mathrm{out}})$ is a completely-positive,
trace-preserving (CPTP) linear map. The Choi-Jamiołkowski theorem
provides an equivalent representation that is central to our
formalism: every CPTP map $\Phi$ is uniquely associated with a
positive semidefinite matrix $J(\Phi) \in \mathcal{B}(\mathcal{H}_{\mathrm{in}}
\otimes \mathcal{H}_{\mathrm{out}})$ called the \textbf{Choi matrix}:

\begin{equation}
J(\Phi) = (\operatorname{id} \otimes \Phi)\!\bigl(
  \ket{\Omega}\!\bra{\Omega} \bigr),
\qquad
\ket{\Omega} = \frac{1}{\sqrt{d_{\mathrm{in}}}}
\sum_{i=0}^{d_{\mathrm{in}}-1} \ket{i} \otimes \ket{i},
\label{eq:choi_def}
\end{equation}
where $d_{\mathrm{in}} = \dim \mathcal{H}_{\mathrm{in}}$ and
$\{\ket{i}\}$ is a fixed orthonormal basis. Conversely, $\Phi$ can be
reconstructed from $J(\Phi)$ via
\begin{equation}
\Phi(\rho) = \operatorname{tr}_{\mathrm{in}}\!\bigl[
  (I_{\mathrm{out}} \otimes \rho^\top) \, J(\Phi) \bigr].
\end{equation}

The Choi matrix admits a convenient \textbf{Kraus representation}: if
$\{K_l\}$ are operators satisfying $\sum_l K_l^\dagger K_l = I$, then
\begin{equation}
J\!\bigl(\Phi\bigr) = \sum_l \ket{K_l}\!\bra{K_l},
\qquad
\ket{K_l} = (I \otimes K_l) \ket{\Omega},
\end{equation}
and the channel action is $\Phi(\rho) = \sum_l K_l \, \rho \, K_l^\dagger$.
This representation is the foundation for the QHMM channel model used in
this work.

A crucial property is \textbf{trace preservation}: $\Phi$ is TP if and
only if $\operatorname{tr}_{\mathrm{out}}[J(\Phi)] = I_{\mathrm{in}}$.
Complete positivity ensures that all eigenvalues of $J(\Phi)$ are
non-negative. The CPTP constraints are essential for the physical
validity of any learned channel.

\index{Choi matrix}
\index{Kraus operators}
\index{CPTP}

% =============================================================================
% §2.4  Hilbert-Schmidt Space and Observable Operator Models
% =============================================================================

\section{Observable Operator Models and the Hilbert-Schmidt Representation}
\label{sec:hmm_oom}

The Observable Operator Model (OOM), introduced by
\citet{jaeger2000oom}, is a generalization of the HMM that represents
the dynamics of a stochastic process as linear operators on a
finite-dimensional vector space. Consider a QHMM with state space
dimension $S$, action set $\mathcal{A} = \{0, \ldots, A-1\}$, and
outcome set $\mathcal{O} = \{0, \ldots, O-1\}$. Let $\mathcal{B}(\mathcal{H}_S)$
be the space of $S \times S$ density matrices, and let
$\{\mathcal{P}_\mu\}_{\mu=0}^{S^2-1}$ be an orthonormal basis of
$\mathcal{B}(\mathcal{H}_S)$ under the Hilbert-Schmidt (HS) inner product
$\langle P, Q \rangle_{\mathrm{HS}} = \operatorname{tr}[P^\dagger Q]$.
A natural choice is the matrix-unit basis
$P_{ij} = \ket{i}\!\bra{j}$.

\begin{definition}[HS Vector]
\index{Hilbert-Schmidt vector}
For a density matrix $\rho \in \mathcal{B}(\mathcal{H}_S)$, its
\textbf{HS vector} is
\begin{equation}
v(\rho) = \bigl( \langle \mathcal{P}_0, \rho \rangle_{\mathrm{HS}},
               \ldots,
               \langle \mathcal{P}_{S^2-1}, \rho \rangle_{\mathrm{HS}} \bigr)^\top
\in \mathbb{R}^{S^2}.
\end{equation}
\end{definition}

The crucial observation is that the dynamics of the QHMM can be encoded
in a set of linear operators on $\mathbb{R}^{S^2}$. Define, for each
transition label $(o, a, o', a')$ (current outcome, current action, next
outcome, next action), an $S^2 \times S^2$ real matrix $A^{(o,a,o',a')}$
with entries
\begin{equation}
A^{(o,a,o',a')}_{\mu,\nu}
  = \operatorname{tr}\!\bigl[ \mathcal{P}_\mu^\dagger \,
    \Phi^{(a')}_{o'} \!\bigl( \mathcal{E}\bigl(
      \Phi^{(a)}_o(\mathcal{P}_\nu) \bigr) \bigr) \bigr],
\label{eq:oom_A_def}
\end{equation}
where $\Phi^{(a)}_o$ is the instrument branch (a CPTP channel) for
action $a$ and outcome $o$, and $\mathcal{E}$ is the channel modeling the
quantum reservoir dynamics.

With these operators, the OOM update rule for a trajectory
$\tau = (a_1, o_1, \ldots, a_T, o_T)$ is
\begin{equation}
v_{t+1} = A^{(o_t, a_t, o_{t+1}, a_{t+1})} \, v_t,
\qquad v_1 = v(\rho_1),
\label{eq:oom_update}
\end{equation}
where $v_t$ is the HS vector of the \emph{filtered} density matrix after
$t$ steps. The one-step probability (marginal likelihood) is
\begin{equation}
Z_t \equiv p(o_t \mid \tau_{1:t-1}) = \mathbf{1}^\top v_t,
\qquad \mathbf{1} = (1, 1, \ldots, 1)^\top,
\label{eq:oom_marginal}
\end{equation}
and the full trajectory probability is
\begin{equation}
p(\tau) = \mathbf{1}^\top v_T = \prod_{t=1}^T Z_t.
\label{eq:oom_trajectory}
\end{equation}

\index{OOM!transition operator}
\index{OOM!forward algorithm}

% =============================================================================
% §2.5  The Forward-Backward Algorithm for OOM
% =============================================================================

\section{Forward-Backward Algorithm for OOM State Smoothing}
\label{sec:forward_backward}

The OOM framework naturally supports a forward-backward algorithm,
analogous to the Baum-Welch algorithm for classical HMMs, for computing
the posterior distribution over hidden states given a full trajectory.

\index{forward-backward algorithm!OOM}
\index{smoothing}

Let $v_t$ denote the unnormalized forward state (HS vector of the
subnormalized filtered state $\tilde\rho_t$) and $Z_t = \mathbf{1}^\top v_t$
the marginal likelihood at step $t$. Define the \textbf{forward message}
as the normalized state
\begin{equation}
\alpha_t = \frac{v_t}{Z_t} \in \mathbb{R}^{S^2},
\qquad \sum_{\mu=0}^{S^2-1} \alpha_t[\mu] = 1.
\label{eq:alpha_def}
\end{equation}
$\alpha_t[\mu]$ represents $p(\rho_t = \mathcal{P}_\mu \mid \tau_{1:t})$,
the posterior probability of the hidden state being the $\mu$-th basis
element given all observations up to time $t$.

The \textbf{backward message} $\beta_t \in \mathbb{R}^{S^2}$ encodes the
likelihood of future observations:
\begin{equation}
\beta_t[\mu] = p\bigl(\tau_{t+1:T} \mid \rho_t = \mathcal{P}_\mu, a_t, o_t\bigr).
\label{eq:beta_def}
\end{equation}
The backward recurrence is
\begin{equation}
\beta_t = \bigl(A^{(o_t, a_t, o_{t+1}, a_{t+1})}\bigr)^\top \beta_{t+1}
           \big/ Z_t,
\qquad \beta_T = \mathbf{1},
\label{eq:beta_recurrence}
\end{equation}
where the division by $Z_t$ corrects for the normalization used in the
forward pass. The key difference from classical HMMs is that the backward
operator uses the \emph{actual next action} from the trajectory, not the
sum over all possible actions.

\index{backward message}
\index{smoothing posterior}

The \textbf{smoothing posterior}—the state posterior conditioned on the
full trajectory—is obtained by combining forward and backward messages:
\begin{equation}
\xi_t \equiv p\!\bigl(\rho_t = \mathcal{P}_\mu \mid \tau_{1:T}\bigr)
        = \frac{\alpha_t[\mu] \, \beta_t[\mu]}
               {\sum_{\nu=0}^{S^2-1} \alpha_t[\nu] \, \beta_t[\nu]}.
\label{eq:smoothing_posterior}
\end{equation}
The denominator $\bar Z_t = \sum_\mu \alpha_t[\mu]\beta_t[\mu]$ equals
$p(o_{t+1:T} \mid \tau_{1:t})$, the likelihood of the future given the
past, and serves as a local model quality indicator.

\index{smoothing posterior}
\index{log-likelihood}

The log-likelihood of the full trajectory follows from
Eqs.\ \eqref{eq:oom_marginal} and \eqref{eq:oom_trajectory}:
\begin{equation}
\log p(\tau) = \sum_{t=1}^T \log Z_t.
\label{eq:log_likelihood}
\end{equation}

A critical implementation detail: the backward recurrence
\eqref{eq:beta_recurrence} uses $A^{(o_t,a_t,o_{t+1},a_{t+1})\top}$, the
matrix transpose (not the conjugate transpose) of the forward operator.
This is because the forward operator $A$ is real symmetric when the
channel is unital, but more generally it is the HS inner product structure
that determines the adjoint in the OOM representation.

% =============================================================================
% §2.6  TD($\lambda$) Eligibility Traces
% =============================================================================

\section{TD($\lambda$) Eligibility Traces for OOM Learning}
\label{sec:td_lambda}

TD($\lambda$) \citep{sutton1988learning, sutton2000review} provides a
mechanism for temporal credit assignment in reinforcement learning. We
adapt the classical TD($\lambda$) framework to the OOM setting by defining
a value function over the forward message $\alpha_t$ and accumulating
eligibility traces over the state space.

\index{TD($\lambda$)}
\index{eligibility traces}

Define the \textbf{value function} as a linear function of the forward
message:
\begin{equation}
V(\alpha_t; w) = \sum_{\mu=0}^{S^2-1} w_\mu \, \alpha_t[\mu] = w^\top \alpha_t,
\qquad w \in \mathbb{R}^{S^2}.
\label{eq:value_function}
\end{equation}
The TD error at step $t$ is
\begin{equation}
\delta_t = r_t + \gamma \, V(\alpha_{t+1}; w) - V(\alpha_t; w),
\label{eq:td_error}
\end{equation}
where $r_t$ is the reward signal and $\gamma \in [0,1]$ is the discount
factor. In the absence of an external reward, we use the
\textbf{surprisal} $r_t = \log Z_t = \log p(o_t \mid \tau_{1:t-1})$ as an
intrinsic reward, which encourages the model to predict likely observations.

\index{surprisal}

\textbf{Forward eligibility traces} accumulate the history of state
visits:
\index{forward trace}
\begin{equation}
e_t = \gamma \lambda_f \, e_{t-1} + \alpha_t,
\qquad e_0 = \alpha_0,
\label{eq:forward_trace}
\end{equation}
where $\lambda_f \in [0,1]$ is the forward trace decay parameter. When
$\lambda_f = 0$, the trace reduces to the immediate state; when
$\lambda_f = 1$, the trace accumulates uniformly over all history.

\index{backward trace}
\textbf{Backward eligibility traces} propagate TD errors backwards through
time:
\begin{equation}
e^b_t = \alpha_t \, |\delta_t| + \gamma \lambda_b \, e^b_{t+1},
\qquad e^b_T = \alpha_T,
\label{eq:backward_trace}
\end{equation}
where $\lambda_b \in [0,1]$ is the backward trace decay and $|\delta_t|$
weights the trace by the magnitude of the TD error, concentrating credit
assignment on timesteps with large prediction errors.

\index{combined trace}
The \textbf{combined eligibility trace} used for parameter updates is the
element-wise product
\index{combined trace}
\begin{equation}
e^{\mathrm{comb}}_t = e_t \odot \tilde e^b_t,
\qquad
\tilde e^b_t = \frac{e^b_t}{\|e^b_t\|_\infty},
\label{eq:combined_trace}
\end{equation}
which balances recency (from the forward trace) with error magnitude
(from the backward trace).

\index{TD($\lambda$)!parameter update}
The OOM transition operator $A^{(o,a,o',a')}$ is updated using the
combined trace:
\index{parameter update}
\begin{equation}
\Delta A^{(o,a,o',a')} = \eta \, \delta_t \, e^{\mathrm{comb}}_t,
\qquad A^{(o,a,o',a')} \leftarrow A^{(o,a,o',a')} + \Delta A^{(o,a,o',a')},
\label{eq:a_update}
\end{equation}
where $\eta > 0$ is the learning rate. This update rule shifts the OOM
dynamics to increase the probability of trajectories with high rewards,
analogously to the policy gradient update in classical RL.

The forward trace \eqref{eq:forward_trace} provides the \emph{forward
view} of TD($\lambda$), which is statistically biased but has low
variance; the backward trace \eqref{eq:backward_trace} provides the
\emph{backward view}, which is unbiased but has high variance. Their
combination follows the strategy established in classical TD($\lambda$)
\citep{sutton2000review}.

\index{TD($\lambda$)!forward view}
\index{TD($\lambda$)!backward view}

% =============================================================================
% §2.7  OMLE: Maximum Likelihood Estimation for QHMM
% =============================================================================

\section{Online Maximum Likelihood Estimation for QHMM}
\label{sec:omle}

The OOM transition operators $\{A^{(o,a,o',a')}\}$ are parameterized by
the underlying QHMM channels and instruments. Given a dataset of
trajectories $\mathcal{D} = \{\tau^i\}_{i=1}^K$, the OMLE procedure
estimates the QHMM parameters by maximizing the total log-likelihood
subject to CPTP constraints.

\index{OMLE}
\index{CPTP constraints}

Let $J(\mathcal{E}_l)$ be the Choi matrix of the $l$-th channel and
$J(\Phi^{(a)}_o)$ the Choi matrix of the instrument branch for action $a$
and outcome $o$. The optimization variables are these Choi matrices. The
constraints are:
\index{channel!CPTP}
\index{instrument!CPTP}
\begin{subequations}
\label{eq:cptp_constraints}
\begin{align}
J(\mathcal{E}_l) &\succeq 0,
& \operatorname{tr}_{\mathrm{out}}[J(\mathcal{E}_l)] &= I_S,
& \forall\, l, \tag{2.7a} \\
J(\Phi^{(a)}_o) &\succeq 0,
& \sum_{o=0}^{O-1} \operatorname{tr}_{\mathrm{out}}[J(\Phi^{(a)}_o)] &= I_S,
& \forall\, a. \tag{2.7b}
\end{align}
\end{subequations}

\index{channel!trace-preserving}
\index{instrument!trace-preserving}

The SDP objective is
\index{SDP!OMLE}
\begin{equation}
\max_{\{J(\mathcal{E}_l),\, J(\Phi^{(a)}_o)\}}
\; \sum_{i=1}^K \sum_{t=1}^{T_i} \log Z^i_t,
\qquad \text{s.t. } \eqref{eq:cptp_constraints},
\label{eq:omle_objective}
\end{equation}
where $Z^i_t$ is the OOM marginal likelihood \eqref{eq:oom_marginal} for
trajectory $i$ at time $t$, computed using the current parameter estimate.
The problem is solved as a semidefinite program (SDP) using either MOSEK
or SCS as the backend solver. After each SDP update, the Choi matrices
are projected back onto the CPTP constraint manifold via the
Jamiołkowski trace-preserving projection \citep{jiang2022structure}.

\index{forward-backward algorithm!role in OMLE}
The forward-backward algorithm (\S\ref{sec:forward_backward}) plays a
dual role: it provides the smooth posterior \eqref{eq:smoothing_posterior}
needed to evaluate the EM auxiliary function, and it provides the
surprisal rewards $r_t = \log Z_t$ used in the TD($\lambda$) updates
(\S\ref{sec:td_lambda}).

% =============================================================================
% §2.8  Optimistic Planning for Prediction
% =============================================================================

\section{Optimistic Planning for Trajectory Prediction}
\label{sec:optimistic_planning}

Given a learned QHMM, prediction requires selecting the action sequence
that maximizes expected cumulative reward. This is a planning problem
in a finite-horizon MDP. We use the \textbf{optimistic planning}
algorithm \citep{buşoniu2020online}, which maintains a set of candidate
models (one per QHMM partition) and selects the model with the highest
expected value at the root.

\index{optimistic planning}
\index{MDP}

For a set of $M$ candidate QHMM environments
$\{\mathcal{M}^{(k)}\}_{k=1}^M$, each with root value estimate
$V^{(k)}_1(\emptyset)$, the algorithm selects
\begin{equation}
k^* = \arg\max_{k \in [M]} V^{(k)}_1(\emptyset),
\qquad
\pi^* = \arg\max_{\pi \in \Pi^{(k^*)}} V^{\pi}_1(\emptyset),
\label{eq:optimistic_plan}
\end{equation}
where $\Pi^{(k)}$ is the policy tree of the $k$-th model. The selected
model's policy tree $\pi^*$ is then used to generate actions for the
next prediction step. This approach efficiently handles the
combinatorial explosion of action sequences by pruning low-value
candidates early.

\index{policy tree}
\index{value function!root}

This completes the mathematical foundation. The next chapter describes
how these components are integrated into the end-to-end QRC-EV system.


% =============================================================================
% Bibliography
% =============================================================================
\bibliographystyle{plain}
\bibliography{refs}
\addcontentsline{toc}{chapter}{Bibliography}

% =============================================================================
% Appendix
% =============================================================================
\appendix
% =============================================================================
% Appendix A — Hyperparameters
% =============================================================================

\chapter{Hyperparameters and Experimental Protocol}
\label{app:hyperparameters}

\section{Quantum Reservoir}

All QRC experiments use the following reservoir configuration unless
otherwise noted:

\begin{table}[htbp]
  \centering
  \caption{Default reservoir hyperparameters.}
  \label{tab:reservoir_hyperparams}
  \begin{tabular}{@{}llc@{}}
    \toprule
    Parameter & Symbol & Value \\
    \midrule
    Number of qubits & $n$ & 8, 16, 20, 30 (ablation) \\
    Input feature dimension & $d_{\mathrm{in}}$ & 8 \\
    Reservoir feature dimension & $d_{\mathrm{res}}$ & 64 \\
    Spectral radius & $\rho(W)$ & 0.9 \\
    Leaking rate & $\alpha$ & 0.3 \\
    Readout regularization & $\lambda_{\mathrm{ridge}}$ & 1.0 \\
    Input scaling & $\gamma_{\mathrm{input}}$ & 0.5 \\
    Depolarizing probability & $p_{\mathrm{depol}}$ & 0.01 \\
    Shot-based measurement & $N_{\mathrm{shots}}$ & 4096 \\
    \bottomrule
  \end{tabular}
\end{table}

\index{hyperparameters!reservoir}
\index{hyperparameters!ablation}

\section{QHMM-OOM}

\index{hyperparameters!QHMM}

\begin{table}[htbp]
  \centering
  \caption{Default QHMM-OOM hyperparameters.}
  \label{tab:qhmm_hyperparams}
  \begin{tabular}{@{}llc@{}}
    \toprule
    Parameter & Symbol & Value \\
    \midrule
    State dimension & $S$ & 2 \\
    Number of actions & $A$ & 2 \\
    Number of outcomes & $O$ & 2 \\
    Number of channels & $L$ & 1 \\
    SDP solver & -- & MOSEK \\
    OMLE max iterations & $\max_{\mathrm{iter}}$ & 100 \\
    OMLE convergence tol. & $\epsilon_{\mathrm{tol}}$ & $10^{-6}$ \\
    Forward trace decay & $\lambda_f$ & 0.8 \\
    Backward trace decay & $\lambda_b$ & 0.8 \\
    Discount factor & $\gamma$ & 0.9 \\
    Learning rate & $\eta$ & 0.01 \\
    \bottomrule
  \end{tabular}
\end{table}

\index{hyperparameters!OMLE}

\section{Datasets}

\index{hyperparameters!datasets}

\begin{table}[htbp]
  \centering
  \caption{Dataset specifications.}
  \label{tab:dataset_specs}
  \begin{tabular}{@{}lcccc@{}}
    \toprule
    Dataset & $N_{\mathrm{train}}$ & $N_{\mathrm{test}}$ &
    Prediction horizon & Source \\
    \midrule
    Sinusoidal & 4000 & 1000 & 1, 5, 10 & Synthetic \\
    Mackey-Glass ($\tau=17$) & 4000 & 1000 & 1, 5, 10 & Synthetic \\
    NARMA-10 & 4000 & 1000 & 1, 5, 10 & Synthetic \\
    Weekly pattern & 4000 & 1000 & 1, 5, 10 & Synthetic \\
    ACN-Data & 3000 & 1000 & 1, 6, 24 & Real EV \\
    Palo Alto EV & 3000 & 1000 & 1, 6, 24 & Real EV \\
    \bottomrule
  \end{tabular}
\end{table}

\index{datasets!EV charging}
\index{datasets!Mackey-Glass}
\index{datasets!NARMA}

\section{Baselines}

\index{baselines}

\begin{itemize}
  \item \textbf{Ridge regression}: $\alpha = 1.0$, MinMaxScaler for input
    and output normalization.
  \item \textbf{ESN}: $N = 200$ units, spectral radius $= 0.9$,
    leaking rate $\alpha = 0.3$, input scaling $= 0.5$.
  \item \textbf{LSTM}: 2 layers, 128 hidden units per layer,
    dropout $= 0.2$, trained with Adam ($\mathrm{lr} = 10^{-3}$),
    early stopping on validation loss.
  \item \textbf{TFT}: Temporal Fusion Transformer with
    $d_{\mathrm{model}} = 128$, $n_{\mathrm{heads}} = 4$,
    $n_{\mathrm{layers}} = 3$, trained with AdamW.
\end{itemize}

\index{LSTM}
\index{Temporal Fusion Transformer}
\index{Echo State Network}
\index{Ridge regression}

% =============================================================================
% Appendix B — Proofs and Derivations
% =============================================================================

\chapter{Proofs and Derivations}
\label{app:proofs}

\section{Derivation of the OOM Transition Operator}
\index{OOM!derivation}

We derive the OOM transition operator $A^{(o,a,o',a')}$ from the
QHMM dynamics. Starting from the filtered density matrix
$\tilde\rho_t = \mathcal{E}\bigl(\Phi^{(a_t)}_{o_t}(\rho_{t-1})\bigr)$,
the HS vector evolves as:
\index{Hilbert-Schmidt!vectorization}
\index{Kraus operators!composition}

\begin{align}
v_{t+1}[\mu]
  &= \langle \mathcal{P}_\mu, \tilde\rho_{t+1} \rangle_{\mathrm{HS}} \\
  &= \Bigl\langle \mathcal{P}_\mu,
      \sum_{o'} \Phi^{(a_{t+1})}_{o'} \!\bigl(
        \mathcal{E}(\Phi^{(a_t)}_{o_t}(\mathcal{P}_\nu)) \bigr)
    \Bigr\rangle_{\mathrm{HS}} \\
  &= \sum_{o'} \operatorname{tr}\!\bigl[
      \mathcal{P}_\mu^\dagger \,
      \Phi^{(a_{t+1})}_{o'} \!\bigl(
        \mathcal{E}(\Phi^{(a_t)}_{o_t}(\mathcal{P}_\nu)) \bigr)
    \bigr] \\
  &= \sum_{o'} A^{(o_t,a_t,o',a_{t+1})}_{\mu,\nu} \, v_t[\nu],
\end{align}
which recovers Eq.\ \eqref{eq:oom_update} of Chapter~\ref{chap:background}.

\index{OOM!derivation}

\section{TD($\lambda$) Convergence}
\index{TD($\lambda$)!convergence}

The combined trace update \eqref{eq:a_update} is a stochastic gradient
descent step on the TD error objective. Under the standard assumptions of
stochastic approximation \citep{sutton2000review}—specifically bounded
rewards, a contracting value function, and a step-size schedule satisfying
$\sum \eta_t = \infty$, $\sum \eta_t^2 < \infty$—the update converges
almost surely to a fixed point of the TD($\lambda$) operator. The
combining of forward and backward traces does not affect the fixed point
(which is determined by the Bellman operator) but does affect the
variance of the learning trajectory.

\index{Bellman operator}
\index{stochastic approximation}

\section{Choi Matrix Trace-Preservation}
\index{Choi matrix!trace preservation}

For a channel $\Phi$ with Kraus operators $\{K_l\}$, the Choi matrix is
\index{trace preservation!Choi matrix}
\begin{equation}
J(\Phi) = \sum_l \ket{K_l}\!\bra{K_l}.
\end{equation}
The partial trace over the output system is
\index{partial trace}
\begin{align}
\operatorname{tr}_{\mathrm{out}}[J(\Phi)]
  &= \sum_{l} \operatorname{tr}_{\mathrm{out}}\bigl[\ket{K_l}\!\bra{K_l}\bigr] \\
  &= \sum_l (\operatorname{tr}_{\mathrm{out}} \ket{K_l}\!\bra{K_l}) \\
  &= \sum_l K_l^\dagger K_l \\
  &= I_{\mathrm{in}},
\end{align}
where the last line uses the Kraus normalization condition
$\sum_l K_l^\dagger K_l = I$. This establishes that the constraint
\eqref{eq:cptp_constraints}(a) is necessary and sufficient for
trace-preservation.

\index{Kraus normalization}
\index{CPTP!necessary and sufficient}

\section{Numerical Stability of the Forward-Backward Algorithm}
\index{forward-backward!numerical stability}

The forward recurrence $v_{t+1} = A^{(o_t,a_t,o_{t+1},a_{t+1})} v_t$ can
produce numerically vanishing or exploding values for long trajectories.
To mitigate this, we normalize at each step:
\index{log-sum-exp trick}
\index{numerical stability}
\begin{equation}
\log Z_t = \log\!\bigl(\mathbf{1}^\top v_t\bigr),
\qquad
\alpha_t = \frac{v_t}{\mathbf{1}^\top v_t},
\qquad
\tilde v_{t+1} = A^{(o_t,a_t,o_{t+1},a_{t+1})} \alpha_t,
\end{equation}
which is implemented in the code using the log-sum-exp trick to avoid
floating-point overflow. The backward pass uses the same normalization
at each step, ensuring that $\beta_t$ remains in a numerically stable
range.

\index{log-sum-exp}
\index{forward-backward!scaling}


\end{document}
